The current implication from leverage, moments, and a lower predictable-QV bound to Gaussianity after division by the random terminal predictable quadratic variation is false for general adaptive martingale transforms.  Here is an explicit certificate inside the complete-block setup.  Let \(K=2\), \(c=(1,-1)/\sqrt2\), take \(B_n\to\infty\) equal blocks of size \(2q_n\) with fixed counts \((q_n,q_n)\), and let every block have a baseline list containing \(q_n\) values \(+1\) and \(q_n\) values \(-1\).  Set all potential outcomes to zero and set the predictable fits in both arms equal to \(-H_{nb}x_{nbi}\).  Use \(H_{nb}=1\) in the first half of the blocks.  If \(X_n\) is the standardized sum of the first-half complete-randomization scores, use \(H_{nb}=1\) in the second half when \(X_n\ge0\), and \(H_{nb}=2\) there when \(X_n<0\).  This rule is pre-block measurable.  The residuals are \(H_{nb}x_{nbi}\), so overlap, every fixed moment envelope, \(q_n/\log(2B_nK)\to\infty\), leverage \(1/B_n\to0\), and \(V_n\ge v_0s_n^2\) all hold.  The within-block standardized hypergeometric scores are independent before the predictable multiplication, and the finite-population Lindeberg theorem gives independent limits \(X,Y\sim\mathcal N(0,1/2)\) for the two half sums.  Yet the proposed studentized statistic converges to
\[
 T=\frac{X+A(X)Y}{\sqrt{\{1+A(X)^2\}/2}},\qquad A(X)=\mathbf 1\{X\ge0\}+2\mathbf 1\{X<0\}.
\]
Since \(E(Y\mid X)=0\),
\[
 E(T)=E\{X\mathbf 1(X\ge0)\}\left(1-\sqrt{2/5}\right)>0,
\]
so \(T\) is not standard normal.  Thus the original theorem's asserted uniform martingale limit is refuted without ever conditioning on the final path.  Stabilization of \(V_n\) at a deterministic design scale is a standard primitive that rules out precisely this endogenous-time-change obstruction; it is weaker than assuming asymptotic normality and allows the martingale CLT to be derived from the already stated fourth-moment and leverage bounds.